
\documentclass{article}
\usepackage{colortbl}
\usepackage{makecell}
\usepackage{multirow}
\usepackage{supertabular}

\begin{document}

\newcounter{utterance}

\centering \large Interaction Transcript for game `cladder', experiment `full\_v1.5\_default', episode 2287 with refineonly{-}Qwen3.5{-}9B.
\vspace{24pt}

{ \footnotesize  \setcounter{utterance}{1}
\setlength{\tabcolsep}{0pt}
\begin{supertabular}{c@{$\;$}|p{.15\linewidth}@{}p{.15\linewidth}p{.15\linewidth}p{.15\linewidth}p{.15\linewidth}p{.15\linewidth}}
    \# & \multicolumn{2}{c}{Player} && \multicolumn{2}{c}{Game Master} \\
    \hline

    \theutterance \stepcounter{utterance}  
    & & & \multicolumn{4}{p{0.6\linewidth}}{
        \cellcolor[rgb]{0.9,0.9,0.9}{
            \makecell[{{p{\linewidth}}}]{
                \texttt{\tiny{[P1$\langle$GM]}}
                \texttt{You are an expert in causal inference. The following question is not a typical commonsense query, but rather a meticulously designed question created by a professor specializing in causal inference, intended to assess the students' mastery of the course content.} \\
\texttt{Imagine a self{-}contained, hypothetical world with only the following conditions, and without any unmentioned factors or causal relationships: Smoking has a direct effect on effort and college admission. Effort has a direct effect on college admission. The overall probability of smoking is 93\%. The probability of nonsmoking and college admission is 3\%. The probability of smoking and college admission is 62\%. Is the chance of college admission smaller when observing smoking?} \\
\texttt{Start your answer with 'Yes' or 'No', followed by additional reasoning or evidence to support your explanation.} \\
            }
        }
    }
     \\ \\

    \theutterance \stepcounter{utterance}  
    & & \multicolumn{4}{p{0.6\linewidth}}{
        \cellcolor[rgb]{0.95,0.95,0.95}{
            \makecell[{{p{\linewidth}}}]{
                \texttt{\tiny{[P1$\leftrightarrow$P]}}
                \texttt{{[}DRAFT{]}} \\
\texttt{No} \\
\\ 
\texttt{**Reasoning:**} \\
\\ 
\texttt{To determine whether the chance of college admission is smaller when observing smoking, we must compare the conditional probability of admission given smoking, \$P(\textbackslash{}text\{Admission\} | \textbackslash{}text\{Smoking\})\$, with the unconditional probability of admission, \$P(\textbackslash{}text\{Admission\})\$.} \\
\\ 
\texttt{Based on the provided data:} \\
\texttt{1.  **Calculate the unconditional probability of admission, \$P(\textbackslash{}text\{Admission\})\$:**} \\
\texttt{    The problem states that the probability of nonsmoking and college admission is 3\% (\$P(\textbackslash{}neg S \textbackslash{}cap A) = 0.03\$) and the probability of smoking and college admission is 62\% (\$P(S \textbackslash{}cap A) = 0.62\$).} \\
\texttt{    Since these are the only two mutually exclusive ways to be admitted (smoking or not smoking), the total probability is the sum:} \\
\texttt{    \$\$P(\textbackslash{}text\{Admission\}) = P(S \textbackslash{}cap A) + P(\textbackslash{}neg S \textbackslash{}cap A) = 0.62 + 0.03 = 0.65\$\$} \\
\\ 
\texttt{2.  **Calculate the conditional probability of admission given smoking, \$P(\textbackslash{}text\{Admission\} | \textbackslash{}text\{Smoking\})\$:**} \\
\texttt{    We are given that the overall probability of smoking is 93\% (\$P(S) = 0.93\$).} \\
\texttt{    Using the definition of conditional probability:} \\
\texttt{    \$\$P(\textbackslash{}text\{Admission\} | \textbackslash{}text\{Smoking\}) = \textbackslash{}frac\{P(S \textbackslash{}cap A)\}\{P(S)\} = \textbackslash{}frac\{0.62\}\{0.93\} \textbackslash{}approx 0.6667\$\$} \\
\\ 
\texttt{3.  **Compare the probabilities:**} \\
\texttt{    *   \$P(\textbackslash{}text\{Admission\} | \textbackslash{}text\{Smoking\}) \textbackslash{}approx 66.7\textbackslash{}\%\$} \\
\texttt{    *   \$P(\textbackslash{}text\{Admission\}) = 65\textbackslash{}\%\$} \\
\\ 
\texttt{Since \$66.7\textbackslash{}\% > 65\textbackslash{}\%\$, the chance of college admission is **larger** (not smaller) when observing smoking.} \\
\\ 
\texttt{*Note on Causal Structure:* While the causal graph (\$\textbackslash{}text\{Smoking\} \textbackslash{}to \textbackslash{}text\{Effort\} \textbackslash{}to \textbackslash{}text\{Admission\}\$ and \$\textbackslash{}text\{Smoking\} \textbackslash{}to \textbackslash{}text\{Admission\}\$) suggests a confounding scenario where smoking might be associated with lower effort (a negative indirect effect), the specific numerical parameters provided in this hypothetical world result in a positive association between smoking and admission. The direct effect of smoking on admission or the distribution of effort among smokers must be sufficiently strong to outweigh any negative indirect effects, resulting in a higher admission rate for smokers in this specific dataset. The question asks strictly about the observed probability based on the given numbers, not the counterfactual causal effect.} \\
            }
        }
    }
    & \\ \\

    \theutterance \stepcounter{utterance}  
    & \multicolumn{4}{p{0.6\linewidth}}{
        \cellcolor[rgb]{0.561,0.737,0.561}{
            \makecell[{{p{\linewidth}}}]{
                \texttt{\tiny{[P1$\rangle$GM]}}
                \texttt{No} \\
\\ 
\texttt{**Reasoning:**} \\
\\ 
\texttt{To determine whether the chance of college admission is smaller when observing smoking, we must compare the conditional probability of admission given smoking, \$P(\textbackslash{}text\{Admission\} | \textbackslash{}text\{Smoking\})\$, with the unconditional probability of admission, \$P(\textbackslash{}text\{Admission\})\$.} \\
\\ 
\texttt{Based on the provided data:} \\
\texttt{1.  **Calculate the unconditional probability of admission, \$P(\textbackslash{}text\{Admission\})\$:**} \\
\texttt{    The problem states that the probability of nonsmoking and college admission is 3\% (\$P(\textbackslash{}neg S \textbackslash{}cap A) = 0.03\$) and the probability of smoking and college admission is 62\% (\$P(S \textbackslash{}cap A) = 0.62\$).} \\
\texttt{    Since these are the only two mutually exclusive ways to be admitted (smoking or not smoking), the total probability is the sum:} \\
\texttt{    \$\$P(\textbackslash{}text\{Admission\}) = P(S \textbackslash{}cap A) + P(\textbackslash{}neg S \textbackslash{}cap A) = 0.62 + 0.03 = 0.65\$\$} \\
\\ 
\texttt{2.  **Calculate the conditional probability of admission given smoking, \$P(\textbackslash{}text\{Admission\} | \textbackslash{}text\{Smoking\})\$:**} \\
\texttt{    We are given that the overall probability of smoking is 93\% (\$P(S) = 0.93\$).} \\
\texttt{    Using the definition of conditional probability:} \\
\texttt{    \$\$P(\textbackslash{}text\{Admission\} | \textbackslash{}text\{Smoking\}) = \textbackslash{}frac\{P(S \textbackslash{}cap A)\}\{P(S)\} = \textbackslash{}frac\{0.62\}\{0.93\} \textbackslash{}approx 0.6667\$\$} \\
\\ 
\texttt{3.  **Compare the probabilities:**} \\
\texttt{    *   \$P(\textbackslash{}text\{Admission\} | \textbackslash{}text\{Smoking\}) \textbackslash{}approx 66.7\textbackslash{}\%\$} \\
\texttt{    *   \$P(\textbackslash{}text\{Admission\}) = 65\textbackslash{}\%\$} \\
\\ 
\texttt{Since \$66.7\textbackslash{}\% > 65\textbackslash{}\%\$, the chance of college admission is **larger** (not smaller) when observing smoking.} \\
\\ 
\texttt{*Note on Causal Structure:* While the causal graph (\$\textbackslash{}text\{Smoking\} \textbackslash{}to \textbackslash{}text\{Effort\} \textbackslash{}to \textbackslash{}text\{Admission\}\$ and \$\textbackslash{}text\{Smoking\} \textbackslash{}to \textbackslash{}text\{Admission\}\$) suggests a confounding scenario where smoking might be associated with lower effort (a negative indirect effect), the specific numerical parameters provided in this hypothetical world result in a positive association between smoking and admission. The direct effect of smoking on admission or the distribution of effort among smokers must be sufficiently strong to outweigh any negative indirect effects, resulting in a higher admission rate for smokers in this specific dataset. The question asks strictly about the observed probability based on the given numbers, not the counterfactual causal effect.} \\
            }
        }
    }
    & & \\ \\

    \theutterance \stepcounter{utterance}  
    & & & \multicolumn{2}{p{0.3\linewidth}}{
        \cellcolor[rgb]{0.9,0.9,0.9}{
            \makecell[{{p{\linewidth}}}]{
                \texttt{\tiny{[GM$|$GM]}}
                \texttt{No} \\
            }
        }
    }
    & & \\ \\

    \theutterance \stepcounter{utterance}  
    & & & \multicolumn{2}{p{0.3\linewidth}}{
        \cellcolor[rgb]{0.9,0.9,0.9}{
            \makecell[{{p{\linewidth}}}]{
                \texttt{\tiny{[GM$|$GM]}}
                \texttt{no} \\
            }
        }
    }
    & & \\ \\

    \theutterance \stepcounter{utterance}  
    & & & \multicolumn{2}{p{0.3\linewidth}}{
        \cellcolor[rgb]{0.9,0.9,0.9}{
            \makecell[{{p{\linewidth}}}]{
                \texttt{\tiny{[GM$|$GM]}}
                \texttt{game\_result = WIN} \\
            }
        }
    }
    & & \\ \\

\end{supertabular}
}

\end{document}
